\section{Introduction}
\label{sec:introduction}

Genetic relatedness between individuals is fundamentally determined by the time to their most recent common ancestor (TMRCA). Under the coalescent~\cite{kingman1982coalescent}, sequence identity between two haplotypes is a direct estimator of this quantity: $I = 1 - 2\mu t$, where $\mu$ is the per-base mutation rate and $t$ is the TMRCA. This relationship is exact to first order under the infinite-sites model and makes pairwise identity a sufficient statistic for TMRCA at each genomic locus. Four independent research programs have recently converged on the conclusion that TMRCA-based relatedness is superior to genotype-based measures: the expected GRM (eGRM) computed from ancestral recombination graphs outperforms the canonical GRM in 97.5\% of simulated association settings~\cite{fan2022egrm}; the Branch GRM derived from genealogical trees provides a theoretically optimal measure of genetic relatedness~\cite{lehmann2026branch_grm}; GeSi classifies relatedness estimators into ``full'' measures (eGRM, Branch GRM) that capture complete genealogical information versus ``shallow'' measures (standard GRM, KING) that lose deep coalescent signal~\cite{gesi2025}; and Super Admixture connects admixture proportions directly to the coancestry matrix $\Theta = Q'\Lambda Q$, where per-population identity means encode the coancestry parameter~\cite{chen2025super_admixture}.

Despite this theoretical consensus, accessing the TMRCA signal has required either computationally expensive ancestral recombination graph (ARG) inference or discrete genotype data processed through variant-calling pipelines. Current state-of-the-art methods for identity-by-descent (IBD) detection---hap-ibd~\cite{browning2020hap_ibd}, IBDseq~\cite{browning2013ibdseq}, GERMLINE~\cite{gusev2009germline}---and local ancestry inference (LAI)---RFMix~\cite{maples2013rfmix}, FLARE~\cite{browning2023flare}---rely on phased genotype data from SNP arrays or short-read sequencing. These approaches introduce three systematic distortions of the underlying TMRCA signal: (i) ascertainment bias from pre-selected marker sets that capture $\sim$85\% of common variation in European populations but only 55--65\% in African populations~\cite{lachance2012ascertainment}; (ii) blindness to structural variants, which contribute $\sim$20~Mb of pairwise differences compared to $\sim$3~Mb from SNPs~\cite{nurk2022t2t}; and (iii) switch errors from statistical phasing that fragment true IBD segments and degrade with distance from heterozygous sites.

The Human Pangenome Reference Consortium (HPRC) has produced high-quality haplotype-resolved assemblies for hundreds of diverse human genomes~\cite{liao2023hprc}. These assemblies provide chromosome-scale phasing resolved by the assembly process itself, and pangenome graphs constructed via minigraph-cactus~\cite{hickey2024minigraph_cactus} capture the full spectrum of genomic variation. Tools such as \texttt{impg}~\cite{impg} enable efficient computation of pairwise sequence identity directly from pangenome alignments. Crucially, pairwise identity from assemblies is a more direct estimator of $1 - 2\mu t$ than any VCF-derived quantity: it captures all variation simultaneously---SNPs, indels, and structural variants---without ascertainment bias, reference bias, or phasing error. The per-window identity matrix across $n$ haplotypes is, in effect, a local expected GRM estimated directly from assemblies, without requiring VCF generation, statistical phasing, or ARG inference.

In this work, we present HPRCv2-IBD (\texttt{impop\_k}), a Rust toolkit for IBD detection, local ancestry inference, and relatedness estimation directly from pangenome alignments. A two-state HMM on windowed sequence identity detects IBD segments, an $N$-state HMM with pairwise contrast emissions infers local ancestry, and Jacquard coefficient estimation enables relatedness analysis---all from the same identity observations. Operating on whole-genome alignments provides several structural advantages: chromosome-scale phasing resolved by the assembly process, eliminating switch errors; capture of all variation simultaneously---SNPs, indels, and structural polymorphisms---without the information loss of variant calling; freedom from marker ascertainment bias, with detection power scaling with coalescent diversity rather than genotyping platform coverage; and a unified framework grounded in coalescent theory.

We validate against hap-ibd and RFMix v2 on chromosomes 10, 11, 12, and 20 using 223 HPRC individuals. Simulation achieves 97.95\% ancestry concordance, exceeding RFMix (95.9\%) by 2.05\,pp; grandparent painting in the Platinum pedigree achieves 99.53\% mean accuracy from assembly alignments alone.
